\section{Five-edge cuts: exact identities and partial restrictions}
\label{sec:five-cut-partial}

This section records exact identities and initial restrictions for the five-edge-cut analysis.
The later universal functional completes the sharp-order cycle bound.

Let $X,Y$ be connected normalized five-poles of ranks $h,k$.
Join corresponding ports by five edges, to obtain $G$ of rank $r=h+k+4$.
For each pair $e=ij$, write $P_e,Q_e$ for the terminal-path counts.
For each four-element port set $S$, and each of its three pairings $\alpha$, write
$f_{S,\alpha},g_{S,\alpha}$ for the counts of two vertex-disjoint paths with that pairing.
Put $F_S=\sum_\alpha f_{S,\alpha}$ and $J_S=\sum_\alpha g_{S,\alpha}$.
The fifth port has even degree in each counted forest.

\begin{proposition}[Five-cut gluing]
\label{prop:five-cut-gluing}
One has the exact identity
\[
 c(G)=c(X)+c(Y)+\sum_{|e|=2}P_eQ_e
 +\sum_{|S|=4}\left(F_SJ_S-\sum_\alpha f_{S,\alpha}g_{S,\alpha}\right).
\]
\end{proposition}
\begin{proof}
A cycle uses zero, two, or four edges of the cut.
The zero-edge and two-edge cases give the first three terms.
For a fixed four-edge set, each shore supplies a pairing forest.
Equal pairings give two cycles, while distinct pairings give one cycle.
Thus the four-edge contribution has the displayed form.
\end{proof}

Normalize all shore counts by $2^h$ or $2^k$ as appropriate.
Write $a_e=2^{-h}P_e$, $u_{S,\alpha}=2^{-h}f_{S,\alpha}$, and $U_S=\sum_\alpha u_{S,\alpha}$.
Every coordinate $a_e$ and every sum $U_S$ lies in $[0,1]$.
Define the shore energy
\[
 E(X)=\sum_e a_e^2+\sum_S U_S^2.
\]
The normalized crossing contribution is at most $(E(X)+E(Y))/2$.
The leading one-half threshold for that contribution is $8$.
This energy estimate drops the nonnegative equal-pairing subtraction, so it can lose important information.

\begin{proposition}[Capping three ports]
\label{prop:five-pole-three-cap}
Choose three ports $A$ and let $d,e$ be the other two.
Add a new vertex adjacent to the three ports of $A$.
The resulting two-pole $N$ has rank $h+2$, and
\[
 P_{de}(N)=P_{de}(X)+\sum_{\{a,b\}\subset A}
 \left(F_{\{d,e,a,b\}}-f_{\{d,e,a,b\},de|ab}\right),
\]
\[
 c(N)=c(X)+\sum_{\{a,b\}\subset A}P_{ab}(X).
\]
For any fixed $\delta,\eta>0$, if some $a_{ab}\ge\delta$ with $a,b\in A$, then, for all sufficiently large $h$,
\[
 a_{de}+\sum_{\{a,b\}\subset A}
 \left(U_{\{d,e,a,b\}}-u_{\{d,e,a,b\},de|ab}\right)\le2+\eta.
\]
\end{proposition}
\begin{proof}
A terminal path avoids the new vertex or traverses two of its incident edges.
In the latter case, deleting the new vertex gives either of the two pairings other than $de|ab$.
A cycle through the new vertex gives an $a$--$b$ path in $X$.
This proves the identities.
If the last inequality failed, $N$ would have normalized terminal-path density greater than $1/2+\eta/4$.
The fixed-density path--cycle tradeoff would give $c(N)=O_\eta(2^{h/2})$.
This contradicts $c(N)\ge\delta2^h$ for large $h$.
\end{proof}

\begin{proposition}[Adding one edge]
\label{prop:five-pole-added-edge}
Choose two ports $i,j$ and let $R$ be the other three.
Adding edge $ij$ gives a tripole $T$ of rank $h+1$.
For each pair $e\subset R$,
\[
 2^{-h}P_e(T)=a_e+U_{\{i,j\}\cup e}-u_{\{i,j\}\cup e,ij|e}.
\]
Consequently,
\[
 \sum_{e\subset R}\left(a_e+U_{\{i,j\}\cup e}-u_{\{i,j\}\cup e,ij|e}\right)
 \le4+O(2^{-h/2}).
\]
If two of the three summands exceed $1+\eta$ for fixed $\eta>0$, the third is $O_\eta(2^{-h/2})$.
\end{proposition}
\begin{proof}
Separate paths that use the added edge from those that avoid it.
Apply the universal tripole path-sum estimate at rank $h+1$.
For the last assertion, divide all three counts by $2^{h+1}$ and apply the two-high-pairs tripole theorem.
\end{proof}

\begin{proposition}[A gap when all path coordinates are positive]
\label{prop:five-pole-positive-energy-gap}
For all sufficiently large $h$, a normalized five-pole with $a_e\ge1/100$ for all ten pairs satisfies
\[
 E(X)\le7.9483<8.
\]
\end{proposition}
\begin{proof}
Call a pair high if $a_e>51/100$.
The two-high-pairs tripole theorem, after suppressing the unused ports, implies that a wedge of high pairs has its third pair $O(2^{-h/2})$.
Thus the high-pair graph is a matching for large $h$, and it has at most two edges.

Apply the four-pole added-edge constraint to each four-port set, with $\delta=1/100$ and $\eta=1/1000$.
Unused degree-two ports may be suppressed.
If suppression creates a loop, that loop is at a nonterminal vertex.
Every path in the added-edge two-pole avoids it.
Its normalized path count is then at most one half, which gives the same constraint directly.
Summing its six oriented constraints gives
\[
 \sum_{e\subset S}a_e+4U_S\le6(1+\eta).
\]
Summing over the five four-port sets counts each pair three times, so
\[
 3\sum_ea_e+4\sum_SU_S\le30(1+\eta).
\]
Since $U_S^2\le U_S$,
\[
 E(X)\le\frac{15}{2}(1+\eta)+\sum_e(a_e^2-3a_e/4).
\]
The summand is at most $1/4$ on a high pair.
For a nonhigh pair it is at most $-0.0074$, by convexity on $[0.01,0.51]$.
There are at most two high pairs.
Therefore
\[
 E(X)\le7.5(1.001)+2(0.25)-8(0.0074)=7.9483.
\]
\end{proof}

The independent script \texttt{check\_five\_cut\_states.py} verifies 1,920 complete five-cut gluings and 40 three-port cappings on four small shores.
The checks include a shore with a bridge.
The files \texttt{explore\_five\_cut\_profiles.py} and \texttt{five-cut-profile-lp-exploration.json} contain numerical relaxations only; they are not proof certificates.
The universal functional proved later resolves both sparse profiles and bounded shores for the cycle bound.
